\chapter{Wstęp i założenia}
\section{Wprowadzenie}
Krajowy System e-Faktur, w skrócie KSeF, to platforma online pozwalająca na wystawianie, odbieranie oraz przechowywanie faktur ustrukturyzowanych, opracowana przez Ministerstwo Finansów \parencite{ministerstwo_finansow_pytania_2025}. 
W trakcie powstawania tej pracy nie występuje wciąż obowiązek wystawiania faktur za pomocą platformy - jednak obowiązek ten będzie stopniowo wprowadzany na przestrzeni 2026 roku \parencite{ministerstwo_finansow_pytania_2025}. 

Pierwszy etap wdrożenia rozpoczyna się 1 lutego 2026 roku, dla firm, których sprzedaż w roku 2024 przekroczyła wartość 200 milionów PLN, wliczając VAT. Drugi, bardziej znaczący dla projektu realizowanego w tej pracy termin przypada na 1 kwietnia 2026 - poczynając od tego dnia, obowiązek wykorzystywania narzędzia obowiązuje wszystkich, z jednym wyjątkiem.

W związku z szeroką gamą zmian, które spowoduje wdrożenie systemu, Ministerstwo Finansów pozostawia możliwość wystawiania faktur poza platformą do 31 grudnia 2026 roku, o ile suma podatku należnego VAT z owych faktur nie przekroczy wartości 10 000 PLN na przestrzeni rozliczanego miesiąca \parencite{ministerstwo_finansow_zakres_2025}. 

\section{Cel i zakres pracy}
Celem pracy jest zaprojektowanie oraz utworzenie oprogramowania umożliwiającego bezpieczną oraz \enquote{bezbolesną} % TODO: znaleźć mniej kolokwialny synonim - "minimalizującą koszty i skomplikowanie"?
integrację z Krajowym Systemem e-Faktur, w szczególności dla małych i średnich przedsiębiorstw. Projekt zrealizowany będzie w języku Go % TODO: przypis? ew. napomnienie o często używanej nazwie Golang
i przybierze formę pojedynczego pliku wykonywalnego, kompatybilnego z systemami Windows, MacOS oraz Linux. 
\par
Zakres pracy obejmuje:
\begin{itemize}
    \item analizę wymagań technicznych oraz oficjalnej dokumentacji Krajowego Systemu e-Faktur
    \item zaprojektowanie architektury oprogramowania,
    \item implementację oprogramowania,
    \item testowanie jego działania,
    \item praktyczne wykorzystanie go w realnym scenariuszu oraz ocenę działania na podstawie wyników.
\end{itemize}
% TODO: po ukończeniu pracy - przypisy do odpowiednich sekcji pracy

\section{Założenia projektu}
Projekt realizowany jest zakładając scenariusz, w którym wiele mniejszych przedsiębiorców, którzy nie wpisują się w kategorie obejmowane warunkami towarzyszącymi terminowi pierwszemu lub drugiemu, zdecyduje się na rozwiązanie problemu poprzez zatrudnienie księgowego, gdy będzie to potrzebne - zleceniowo w przypadku przekroczenia progu podatku należnego w rozliczanym miesiącu, lub na stałe od początku roku 2027 \parencite{ministerstwo_finansow_zakres_2025}. 

Oprogramowanie tak więc kierowane jest wobec mniejszych przedsiębiorców, dla których integracja z Krajowym Systemem e-Faktur może wiązać się z komplikacjami oraz znaczącym podniesieniem kosztów - i na podstawie tego założenia zostaną przyjęte wszystkie inne.
\par
Założenia:
\begin{itemize}
    \item oprogramowanie kierowane jest wobec małych i średnich przedsiębiorstw (bis)
    \item oprogramowanie może być dostosowane do wymagań przedsiębiorstwa % TODO: w którejś części pracy zamieścić informacje o możliwości integracji z systemami ERP, Excelem, etc.
    \item oprogramowanie nie wymaga wiedzy na temat funkcjonalności Krajowego Systemu e-Faktur
    \item oprogramowanie jest kompatybilne z trzema najpopularniejszymi systemami operacyjnymi w Polsce, na dzień przygotowywania tej pracy \parencite{statcounter_desktop_2025}
    \item oprogramowanie nie wymaga nowoczesnego sprzętu komputerowego % TODO: w którejś części pracy zamieścić wymagania sprzętowe
    \item oprogramowanie przewiduje możliwe zakłócenia w komunikacji z Krajowym Systemem e-Faktur - przez niedostępność serwisu lub problemy z połączeniem internetowym - oraz zapewnia odpowiednie mechanizmy bezpieczeństwa do radzenia sobie z nimi
    \item oprogramowanie nie ingeruje w treść faktur i zapewnia, że dane przekazywane do Urzędu Skarbowego poprzez KSeF odzwierciedlają dane wprowadzone przez użytkownika.
    %TODO: dodatkowe wymagania, wspomnieć o panelu zatwierdzania?
\end{itemize}

\section{Ogólna architektura oprogramowania}
Oprogramowanie, którego dotyczy ta praca, pełni rolę pośrednika między przedsiębiorcą lub pracownikiem a Krajowym Systemem e-Faktur. Jego główną rolą jest przyjęcie faktur w wybrany przez użytkownika sposób, przekonwertowanie ich do formatu akcepotwanego przez KSeF, wysyłka oraz upewnienie się, że proces został wykonany bez żadnych nieprawidłowości - lub poinformowanie użytkownika o działaniach wymaganych do poprawnego wykonania procesu.
\par
Na oprogramowanie składają się następujące elementy:
\begin{itemize}
    \item interfejs wejściowy, przyjmujący dane bezpośrednio od przedsiębiorcy lub systemów zewnętrznych
    \item oprogramowanie przetwarzające faktury i konwertujące je do formatu akceptowanego przez KSeF
    \item oprogramowanie komunikujące się z Krajowym Systemem e-Faktur, odpowiedzialne za wysyłanie faktur oraz pobieranie informacji na temat ich stanu - m.in. Urzędowe Poświadczenia Odbioru
    \item lokalna baza danych do przechowywania danych na temat statusu wysyłki poszczególnych faktur
    \item \enquote{kolejka} lokalna, do przechowywania danych w przypadku przerwania komunikacji z internetem/braku dostępności KSeF 
    \item panel z interfejsem graficznym przeznaczonym dla użytkownika, informujący o statusie wysyłki lub/oraz o napotkanych nieprawidłowościach
\end{itemize}

Oprogramowanie przyjmuje faktury od przedsiębiorcy/systemu zewnętrznego oraz przetwarza dane, przygotowywując je do wysyłki i dodając do lokalnej kolejki. Następnie, warstwa odpowiedzialna za komunikację z KSeF weryfikuje połączenie, wysyła dane i nasłuchuje odpowiedzi, zależnie od której aktualizuje status faktury. W przypadku niepowodzenia, oprogramowanie dokonuje dodatkowych prób i ostatecznie oznacza status wysyłki jako niepowodzenie.

\section{Narzędzia wykorzystane do realizacji}
\subsection{Język programowania Go}
Go, znany również jako \enquote{Golang}, jest językiem programowania open-source o nieskomplikowanej składni, wyposażonym w rozległą bibliotekę standardową i zbudowanym z myślą o współbieżności. Pomimo podobieństwa składni do C, Go automatycznie zarządza pamięcią. Jest również językiem kompilowanym i wspiera wiele platform, a w branży informatycznej stosowany jest głównie do serwisów back-end. Biorąc pod uwagę wyżej wymienione cechy oraz ich dopasowanie do założeń projektu, został on wybrany do realizacji oprogramowania. Na dzień przygotowywania tej pracy, najnowszą dostępną wersją Go jest 1.25.5 - w niej również zrealizwany będzie projekt. % TODO: przypis

\subsection{Baza danych SQLite}
SQLite jest lekką, relacyjną bazą danych, która nie wymaga dedykowanego serwera. Wszystkie dane przechowywane są w jednym pliku, lokalnie. Jest to przenośne rozwiązanie, które ułatwi implementację warstwy aplikacji odpowiedzialnej za działanie w razie zakłóceń z połączeniem internetowym/dostępnością KSeF. 

\subsection{Biblioteki/pakiety}
Projekt będzie korzystał przede wszystkim z biblioteki standardowej Go - w tym pakietów takich jak net, crypto, encoding i database. % TODO: przypis
Dodatkowo, by zapewnić integrację z SQLite, wykorzystany będzie pakiet modernc.org/sqlite % TODO: przypis
jest on napisany całkowicie w Go, wpisując się w założenie o przenośności.

% TODO: docker
% TODO: dopytać, czy powinienem również opisać środowisko, w którym oprogramowanie było tworzone
% TODO: wspominać o gicie?
% TODO: zdecydować się, czy określać ostatecznego odbiorcę mianem "przedsiębiorcy" czy "użytkownika" - w tekście występują oba
